\section{Related Work}
\label{sec:related_work}

\para{Superposition and the dimensionality bottleneck.}
\citet{elhage2022superposition} established that neural networks store more features than they have dimensions by packing sparse features as nearly-orthogonal directions---the superposition hypothesis.
\citet{scherlis2022polysemanticity} showed that polysemanticity emerges naturally from capacity constraints.
These results imply that compound concepts, being relatively sparse in training data compared to their constituent words, are likely stored in superposition with related concepts.
Our work tests a more specific question: whether compound concepts have \emph{any} dedicated direction (even in superposition) or are entirely derived from constituent directions augmented by next-token prediction.

\para{Linear representations in LLMs.}
The linear representation hypothesis posits that concepts correspond to directions in a model's activation space.
\citet{park2024linear} formalized three variants and tested 27 single-token concepts in Llama-2, finding strong linear structure for most.
\citet{gurnee2023language} showed that LLMs develop linear representations of space and time.
\citet{merullo2023language} demonstrated that LLMs implement Word2Vec-style vector arithmetic for relational tasks, suggesting that compositional operations are possible within the residual stream.
Unlike these studies, which focus on single-token concepts or binary relations, we directly test whether \emph{multi-token} compound concepts have independent directions or are compositionally derived from their parts.

\para{Sparse autoencoders and feature decomposition.}
Sparse autoencoders (SAEs) decompose polysemantic activations into more interpretable features~\citep{cunningham2023sparse,gao2024scaling,rajamanoharan2024gated}.
However, \citet{chanin2024absorption} identified the feature absorption problem: when features form hierarchies, SAEs absorb parent features into child features, creating unreliable classifiers.
\citet{minegishi2025rethinking} showed that SAEs can separate different word senses for polysemous words, with sense disambiguation concentrated in deeper layers.
\citet{giglemiani2024evaluating} demonstrated that arbitrary combinations of SAE latents do not produce realistic activations, suggesting significant geometric structure beyond simple addition.
We use SAE feature analysis not as ground truth but as a complementary lens: the finding that 56\% of compound features are unique provides evidence of compound-specific computation that cosine similarity misses, while acknowledging that SAE artifacts may affect exact percentages.

\para{Compound noun semantics in transformers.}
\citet{ormerod2024kitchen} used representational similarity analysis to compare transformer representations of compound nouns with human judgments in BERT-family models, finding that compounds processed together have different representations than constituents processed separately.
\citet{miletic2023systematic} tested 41,496 BERT configuration combinations for predicting compound compositionality, finding that early layers and surrounding context are most informative.
\citet{nedumpozhimana2024semantics} surveyed multiword expression processing in transformers, confirming that these models encode MWE semantics but with heavy dependence on extraction method and layer choice.
Building on these findings, we extend compound analysis to autoregressive models and introduce SAE feature decomposition and next-token priming as new analytical tools, providing a more fine-grained picture than residual stream geometry alone.

\para{Combinatorial interpretability.}
\citet{adler2025combinatorial} proposed feature channel coding as an alternative to SAE-based analysis, where compound features are encoded via overlapping neuron groups in MLP weight matrices.
This framework predicts that compound concepts should be represented through intersecting feature channels rather than dedicated directions---consistent with our finding that compounds are contextual modifications of component representations rather than independent directions.
